\documentclass[a4paper,11pt]{article}
\usepackage{exam-style-341}
\examinehead{341 农业综合知识三 · 网络技术分章练习}{第07章 运输层 参考答案与解析}
\begin{document}

\unitanswerscover{341 农业综合知识三 · 网络技术分章练习}{第07章 运输层 参考答案}{第07章练习的答案与解析；判断题均给出改正后的表述，序号与确认号、滑动窗口、拥塞窗口计算均给出完整推导，便于复算核对。}

\answerpart{一、单项选择题（10 题，每题 3 分，共 30 分）}

\q C。\begin{analysis}运输层为应用进程之间提供端到端的逻辑通信，用 16 位端口号标识主机中的应用进程；IP 地址标识主机，MAC 地址标识网卡，域名是主机名字符串。IP 地址＋端口号构成套接字。（教材 5.1.1、5.1.3，印刷页 211、214）\end{analysis}

\q B。\begin{analysis}UDP 用户数据报首部只有 4 个字段：源端口、目的端口、长度、检验和，共 8 字节（教材 5.2.2，印刷页 217）。A 4 字节太小；C 20 字节是 TCP 首部固定部分；D 40 字节是 IPv6 基本首部。\end{analysis}

\q B。\begin{analysis}TCP 报文段首部的前一部分是固定部分，共 20 字节；加上选项字段后最长可达 60 字节（教材 5.5，印刷页 225）。A 8 字节是 UDP 首部；C、D 分别是 IPv6 基本首部和 TCP 首部含选项的最大长度。\end{analysis}

\q A。\begin{analysis}端口号 0—1023 为熟知端口（系统端口），1024—49151 为登记端口，49152—65535 为短暂端口。80 是 HTTP 的熟知端口；1024 属于登记端口；50000、65535 属于短暂端口（教材 5.1.3，印刷页 214）。\end{analysis}

\q A。\begin{analysis}三次握手：第一次 A 发送 SYN=1、seq=x；第二次 B 发送 SYN=1、ACK=1、seq=y、ack=x+1；第三次 A 发送 ACK=1、seq=x+1、ack=y+1。因此第二个报文段 SYN 和 ACK 同时置 1（教材 5.9.1，印刷页 247）。\end{analysis}

\q C。\begin{analysis}主动关闭方发送最后一个 ACK 后进入 TIME\_WAIT 状态，等待 2MSL（最长报文段寿命的两倍）后才进入 CLOSED（教材 5.9.2，印刷页 248—249）。A CLOSE\_WAIT 是被动关闭方收到 FIN 后的状态；B LAST\_ACK 是被动关闭方最后确认的状态；D FIN\_WAIT\_1 是主动关闭方发送 FIN 后的状态。\end{analysis}

\q B。\begin{analysis}TCP 发送窗口取接收窗口与拥塞窗口中的较小者，即 $W=\min(rwnd,cwnd)$：既不能超过接收方的缓冲能力（流量控制），也不能超过网络的承载能力（拥塞控制）。A 取大值会同时突破两个限制；C、D 只考虑一个因素。\end{analysis}

\q B。\begin{analysis}慢开始阶段，每经过一个 RTT，cwnd 加倍（1、2、4、8……），呈指数增长；当 cwnd 达到慢开始门限 ssthresh 后转入拥塞避免，每 RTT 加 1（教材 5.8.2，印刷页 241—242）。\end{analysis}

\q A。\begin{analysis}快重传要求接收方每收到一个失序报文段就立即发送重复确认；发送方一旦连续收到 3 个重复确认，就判定该报文段丢失并立即重传，而不必等超时计时器到期（教材 5.8.2）。B 超时触发的是超时重传，走慢开始；C 1 个重复确认不足以判定丢失；D 接收窗口为 0 触发零窗口探测。\end{analysis}

\q B。\begin{analysis}确认号表示期望收到的下一个字节的序号：若确认号为 N，说明序号 N−1 及其之前的字节已全部正确收到。C 是序号的含义；A、D 都不是确认号的含义（教材 5.5）。\end{analysis}

\answerpart{二、多项选择题（3 题，每题 4 分，共 12 分）}

\q ABC。\begin{analysis}TCP 是面向连接的、可靠的、面向字节流的传输协议，支持全双工通信。D 错误：无连接是 UDP 的特点，TCP 通信前必须先建立连接。\end{analysis}

\q ABC。\begin{analysis}UDP 无连接、尽最大努力交付、面向报文、首部仅 8 字节、开销小，并支持一对一、一对多、多对一和多对多通信。D 错误：UDP 不保证可靠交付，可靠性需由应用层负责。\end{analysis}

\q ABC。\begin{analysis}TCP 拥塞控制包括慢开始、拥塞避免、快重传、快恢复四种方法。D 错误：流量控制是接收方通过 rwnd 限制发送方速率，属于端到端的流量控制，不属于拥塞控制方法（二者常配合使用，但概念不同）。\end{analysis}

\answerpart{三、判断题（6 题，每题 2 分，共 12 分）}

\q √。\begin{analysis}TCP 首部固定部分 20 字节（最长 60 字节），UDP 首部 8 字节，这是两者开销差异的重要体现（教材 5.2.2、5.5）。\end{analysis}

\q ×。改正：确认号为 500 表示期望收到的下一个字节的序号是 500，即序号 1—499 的字节已正确收到，而不是“已经收到了 500 字节”。\begin{analysis}确认号是“期望的下一个字节序号”，不是累计已收字节数；若初始序号为 1，收到 499 字节后确认号才是 500。\end{analysis}

\q √。\begin{analysis}若客户端的一个失效 SYN 报文段滞留在网络中，延迟到达服务器，三次握手可使服务器通过第三次 ACK 确认客户端确实想建立连接，避免因失效请求而建立错误连接、浪费资源（教材 5.9.1）。\end{analysis}

\q ×。改正：拥塞窗口 cwnd 并非越大越好，它应与网络承载能力相适应；cwnd 过大会导致网络拥塞、丢包和时延增大，拥塞控制正是要把 cwnd 控制在合理范围。\begin{analysis}发送窗口取 min(rwnd,cwnd)，cwnd 只是上限之一；无限增大 cwnd 会加重网络负担。\end{analysis}

\q √。\begin{analysis}流量控制是端到端的：接收方根据自身缓冲区大小在首部窗口字段中通告 rwnd，发送方据此调整发送窗口，防止接收方缓冲区溢出（教材 5.7.1，印刷页 236）。\end{analysis}

\q √。\begin{analysis}TCP 是全双工协议，连接的两端各自独立发送数据。一方发送 FIN 仅表示自己不再发送数据，但仍可接收；另一方可能还有数据要发，因此 FIN 与 ACK 通常不能合并，需要四次挥手（教材 5.9.2）。\end{analysis}

\answerpart{四、计算题（1 题，共 6 分）}

\q 参考答案与推导：
（1）三次握手中的序号（2 分）：A 的 SYN 报文段序号为 ISN=1000，SYN 占用 1 个序号，故 A 下一个报文段序号为 1001。B 的 SYN+ACK 报文段序号为 ISN=800，确认号为 $1000+1=1001$（确认已收到 A 的 SYN）。A 第三次握手发送 ACK=1、序号 1001、确认号 $800+1=801$；握手完成后 A 下一次发送数据的序号为 1001。
（2）A 发送序号 1001、数据 300 字节（1 分）：本报文段占用序号 1001—1300，B 的确认号 $=1001+300=1301$。
（3）A 发送序号 1301、数据 500 字节（2 分）：本报文段占用序号 1301—1800，B 的确认号 $=1301+500=1801$。
（4）A 发送 FIN、序号 1801（1 分）：FIN 占用 1 个序号（1801），B 的确认号 $=1801+1=1802$。
\begin{analysis}规则：序号指出本报文段数据第一个字节的编号；确认号＝期望收到的下一个字节编号；SYN 和 FIN 虽不携带数据，但各占 1 个序号。因此“下一个序号＝本段起始序号＋数据长度（＋SYN／FIN 的 1）”。\end{analysis}

\answerpart{五、简答题（2 题，每题 10 分，共 20 分）}

\q 参考答案要点（每方面 2 分，共 10 分）：
（1）连接：TCP 面向连接，通信前需三次握手建立连接、通信后四次挥手释放；UDP 无连接，直接发送数据报。
（2）可靠性：TCP 提供可靠交付（序号、确认、重传、校验），保证不丢失、不重复、按序到达；UDP 尽最大努力交付，不保证可靠、不保证按序。
（3）首部与开销：TCP 首部固定 20 字节，含序号、确认号、窗口等大量控制字段；UDP 首部仅 8 字节，开销小。
（4）流量控制与拥塞控制：TCP 有滑动窗口流量控制和慢开始、拥塞避免、快重传、快恢复等拥塞控制；UDP 都没有。
（5）适用场合：TCP 适合对可靠性要求高的应用（HTTP、FTP、SMTP、POP3）；UDP 适合实时性要求高、能容忍少量丢失或需广播／多播的应用（DNS、DHCP、实时音视频）。

\q 参考答案要点：
（1）三次握手过程（3 分）：A 发送 SYN=1、seq=x；B 回复 SYN=1、ACK=1、seq=y、ack=x+1；A 再发送 ACK=1、seq=x+1、ack=y+1，连接建立。
（2）为什么三次（2 分）：两次握手只能确认 A 的发送能力和 B 的接收能力，无法让 A 确认 B 的发送能力和自己的接收能力；三次握手还可防止失效的连接请求报文段突然到达服务器而建立错误连接。
（3）四次挥手过程（3 分）：主动关闭方发送 FIN、seq=u；对方回复 ACK、ack=u+1；对方发送完剩余数据后发送 FIN、seq=w；主动关闭方回复 ACK、ack=w+1，连接释放。
（4）为什么四次（1 分）：TCP 是全双工的，一方发送 FIN 只表示自己不再发送数据，另一方可能还有数据要发送，故收到 FIN 后先回 ACK，待数据发完再发自己的 FIN。
（5）TIME\_WAIT 的作用（1 分）：主动关闭方等待 2MSL，一是保证最后一个 ACK 能到达对方（若丢失可重传），二是使本连接中滞留的旧报文段在网络中消失，避免影响新连接（教材 5.9.1、5.9.2）。

\answerpart{六、综合题（1 题，共 20 分）}

\q 参考答案与推导：
\textbf{（1）第 1—8 个 RTT 的 cwnd（5 分）}
初始 cwnd=1 MSS，ssthresh=16 MSS。
慢开始阶段（cwnd $<$ ssthresh 时每 RTT 加倍）：第 1 RTT 结束 cwnd=1，第 2 RTT=2，第 3 RTT=4，第 4 RTT=8，第 5 RTT=16（达到 ssthresh）。
拥塞避免阶段（cwnd $\ge$ ssthresh 后每 RTT 加 1）：第 6 RTT=17，第 7 RTT=18，第 8 RTT=19。
即：1、2、4、8、16、17、18、19（MSS）。

\textbf{（2）超时后的 ssthresh、cwnd 及之后 5 个 RTT（5 分）}
第 5 RTT 结束后 cwnd=16 MSS 时发生超时：ssthresh $=16/2=8$ MSS，cwnd 重置为 1 MSS（2 分）。
之后：第 6 RTT=1，第 7 RTT=2，第 8 RTT=4，第 9 RTT=8（达到新的 ssthresh=8），随后转入拥塞避免，第 10 RTT=9（3 分）。
即：1、2、4、8、9（MSS）。

\textbf{（3）发送窗口与可发送字节数（4 分）}
发送窗口 $W=\min(rwnd,cwnd)=\min(3000,4000)=3000$ 字节（2 分）。
已发送但未确认 1000 字节，故还能发送 $3000-1000=2000$ 字节，即 2 个 MSS（2 分）。

\textbf{（4）快重传与快恢复（6 分）}
快重传（3 分）：接收方每收到一个失序报文段就立即发送重复确认；发送方若连续收到 3 个对同一报文段的重复确认，就判定该报文段已丢失，立即重传该报文段，而不必等待超时计时器到期，从而缩短恢复时间、提高吞吐量。
快恢复（3 分）：与快重传配合使用。发送方执行快重传后，把 ssthresh 调整为当前 cwnd 的一半，并把 cwnd 设为 ssthresh（或 ssthresh＋3 MSS），不再回到 cwnd=1 的慢开始，而是直接进入拥塞避免，使发送速率不至于下降过多。
\begin{analysis}复核要点：慢开始“加倍”只到 ssthresh 为止；超时后 ssthresh 减半、cwnd 归 1，重新慢开始；快重传由 3 个重复确认触发，快恢复不经过 cwnd=1，这是“超时”与“快重传”两种情形的关键区别。\end{analysis}

\end{document}
